%==============================================================================
% PART 4: Discussion, Conclusion, References
%==============================================================================

\section{Discussion}

\subsection{Key Findings}

This work demonstrates that automated CCD+BAP generation is feasible and accurate for AF3 glycan modeling. The PDB validation provides independent confirmation of stereochemistry correctness:

\begin{enumerate}
\item \textbf{Sialic Acid Handling}: The critical C2 anomeric position for Neu5Ac was correctly identified in all test cases, validating our ketose-specific logic.

\item \textbf{Fucose L-Configuration}: Fucose was correctly mapped to FUC (alpha-L) rather than BDF (beta-L), confirming proper absolute configuration handling.

\item \textbf{Branch Topology}: Multi-residue branched structures were correctly processed, with all CCD codes matching expected values.

\item \textbf{Epimer Distinction}: Galactose vs glucose distinction was maintained through CCD code assignment.
\end{enumerate}

\subsection{Limitations}

Several limitations should be acknowledged:

\begin{enumerate}
\item \textbf{CCD Coverage}: While 28+ monosaccharides are supported, rare sugars may require manual CCD lookup.

\item \textbf{PDB Validation Scope}: Due to AF3 model parameter access restrictions, structure prediction validation was limited to CCD mapping comparison with known PDB structures (n=4). While these four cases were specifically selected to test critical stereochemistry features, the sample size limits generalization.

\item \textbf{Complex Modifications}: Unusual modifications (phosphorylation, sulfation) may require additional handling.

\item \textbf{Processing Time Estimate}: The "30-60 minutes manual BAP specification" claim is based on literature estimates and expert surveys, not direct experimental measurement in this study.
\end{enumerate}

\subsection{Comparison with Related Work}

\begin{table}[h]
\centering
\caption{Comparison with Existing Methods}
\begin{tabular}{lcccc}
\toprule
Method & Stereochemistry & Time & AF3 Ready & Validation \\
\midrule
SMILES & $\sim$60\% & N/A & Yes & Huang et al. 2025 \\
userCCD & $\sim$80\% & N/A & Yes & Huang et al. 2025 \\
Manual BAP & $>$98\% & 30-60 min & Yes & Literature estimate \\
GlycoSMILES2BAP & 97.8\% & $<$1 s & Yes & PDB + Benchmark \\
\bottomrule
\end{tabular}
\label{tab:comparison}
\end{table}

\section{Conclusions}

GlycoSMILES2BAP provides an automated, accurate solution for AF3 glycan input generation. Key contributions include:

\begin{enumerate}
\item \textbf{Automated Pipeline}: Converts IUPAC-condensed and WURCS formats to AF3-compatible CCD+BAP format in less than 1 second per structure.

\item \textbf{Stereochemistry Preservation}: Achieves 97.8\% overall accuracy on benchmark (50 structures) with 100\% CCD mapping accuracy on 4 selected critical PDB test cases.

\item \textbf{Critical Case Validation}: Sialic acid C2 anomeric position correctly handled, confirmed by PDB structure comparison.

\item \textbf{Database Integration}: Compatible with GlyTouCan and GlyGen databases for large-scale processing.
\end{enumerate}

The tool eliminates the input format barrier for AF3 glycan modeling, making accurate structure prediction accessible to the glycobiology community.

\section*{Acknowledgments}

The author thanks the glycobiology community for valuable feedback and the GlyTouCan consortium for maintaining the glycan structure repository.

\section*{Funding}

This work was supported by Zhejiang Xinghe Tea Technology Co., Ltd.

\section*{References}

\bibliographystyle{abbrvnat}
\bibliography{references_final}

\end{document}
